\documentclass[a4paper]{article}
\usepackage{graphicx} % Required for inserting images
\usepackage{wrapfig}
\usepackage{enumitem}
\usepackage[center]{caption}
\usepackage[english, russian]{babel}
\usepackage{amsmath}
\usepackage{geometry}

\geometry{verbose, a4paper, tmargin=2cm, bmargin=2cm, lmargin=2.0cm, rmargin=2.0cm}

\begin{document}

\begin{titlepage}
	\centering
	{\scshape\Large МОСКОВСКИЙ ФИЗИКО-ТЕХНИЧЕСКИЙ ИНСТИТУТ \\
	(НАЦИОНАЛЬНЫЙ ИССЛЕДОВАТЕЛЬСКИЙ УНИВЕРСИТЕТ)\\ % название ВУЗа большим шрифтом
    Физтех-школа аэрокосмических технологий}
	
	\vspace{4cm} % отступ 4 см по вертикали
	{\LARGE Отчет о выполнении лабораторной работы 1.4.2}
	
	\vspace{1cm} % ещё 1 см
	{\huge\bf  Определение ускорения свободного падения при помощи оборотного маятника	}
	
	\vspace{1cm} % ещё 1 см
	\vfill % заполнение по вертикали пустотой (LaTeX сам решит, сколько здесь отступить)
	
    \begin{flushright} % Этот текст будет с правого края страницы
	   {\LARGE Ефремова Татьяна, Б03-503}
    \end{flushright}
    \vfill

	\today % Дата сборки документа
\end{titlepage}

\section{Аннотация}
Цели работы: измерить величину ускорения свободного падения при помощи оборотного маятника и прямых измерений падающего тела.




\section{Теоретические сведения}

\subsection{Теоретическая справка для оборотного маятника}
Физическим маятником называют твердое тело, способное совершать колебания в вертикальной влоскости, будучи подвешенным за одну из своих точек.
При малых колебаниях период колебания физического маятника определяется формулой
\begin{equation} T=2\pi\sqrt{\frac{J}{mgl}}, \end{equation}
где $J$ - момент инерции маятника относительно оси качения, $l$ - расстояние от оси качения до центра масс, $m$ - масса маятника. \\ \\
Приведенная длина маятника - такая длина нити подвеса для математического маятника, что его период совпадает с периодом колебания физического маятника:
\begin{equation}l_{пр}=\frac{J}{ml} \end{equation}
\begin{wrapfigure}{R}{0.2\textwidth}
    \includegraphics[width=1\linewidth]{C:/Users/User/Documents/mipt_takonef/лабы/1.4.2/1.4.2_scheme.png}
    \caption{к т. Гюйгенса}
\end{wrapfigure}

\noindent 
Теорема Гюйгенса об оборотном маятнике: \\
Пусть $O_1$- точка подвеса физического маятника, $C$ -- его центр масс. Отложим отрезок равный $l_{пр}$ на прямой $O_1C$ от точки $O_1$ в направлении $C$. Его конец назовем $O_2$. Точки обладают свойством взаимности: если перевернуть маятник и подвесить его за точку $O_2$, то период его малых колебаний останется таким же, как при подвешивании за точку $O_1$. \\
Пусть $L=l_1+l_2$ - расстояние между двумя точками подвеса физического маятника, лежащими на одной прямой с его центром масс. Если соответствующие периоды малых колебаний равны $T_1 = T_2 = T$, то по теореме Гюйгенса. $L=l_{пр}$. Тогда из (1), (2):
 \begin{equation}g_0=(2\pi)^2\frac{L}{T^2} \end{equation}
Т. к. в реальности невозможно добиться точного совпадения периодов колебаний, из формулы (1) и теоремы Гюйгенса-Штейнера:
 \begin{equation}g=(2\pi)^2\frac{l_1^2-l_2^2}{T_1^2l_1-T_2^2l_2}. \end{equation}
\noindent
Введем $\lambda = \frac{l_1}{l_2}$.
Оценим величину поправки. Пусть $\varepsilon=\frac{|T_2-T_1|}{T_1} \ll 1$ - относительное отклонение при измерении периодов. Тогда при $\lambda\neq1$, пользуясь малостью $\varepsilon$, получим: 
 \begin{equation}g=g_0\frac{\lambda-1}{\lambda-(1+\varepsilon)^2}\approx g_0 \frac{1}{1-\frac{2\varepsilon}{\lambda-1}}\approx g_0(1+2\beta\varepsilon), \end{equation}
где $\beta=\frac{1}{\lambda-1}=\frac{l_2}{l_1-l_2}$. \\ \\
 При $l_2 \rightarrow l_1$, поправка $\delta g$ неограниченно возрастает. Потому будем стремиться выполнять условие: \begin{equation} l_1 > 2,5l_2. \end{equation}

\subsection{Экспериментальная установка оборотного маятника}
Применяемый в работе маятник представляет собой стержень длиной $\approx 1$ м и массой $\approx 1-1.5$ кг. Маятник подвешивается на небольшие треугольные призмы (П1 и П2), острым основанием опираясь на консоль. На стержне закрепляются два дополнительных груза (Г1 и Г2). Для выполнения условия (6), груз Г2 следует закрепить за призмой П2, а Г1 между призмами. \\
Регистрация времени колебаний проводится с помощью электронных счетчиков.  Призмы крепятся на равных расстояних от концов стержня так, чтобы при колебаниях маятник не задевал оправу счетчика.\\
Фиксированное положение призм однозначно задает приведенную длину оборотного маятника $l_{пр} = L$. Главная задача опыта - подобрать такое положение грузов, при котором периоды колебаний при перевороте маятника совпадали бы с достаточно высокой точностью.

\subsection{Теоретический подбор положения грузов для оборотного маятника}
Пусть $b_1$ - расстояние от груза Г1 до призмы П2, а $b_2$ - расстояние от груза Г2 до призмы П2. Масса стержня - $m_\text{ст}$, грузов - $m_1$ и $m_2$. Будем считать положение центра масс заданнымю. \\
{\centering
\begin{figure}[h]
   \centering    
    \includegraphics[width=0.5\linewidth]{C:/Users/User/Documents/mipt_takonef/лабы/1.4.2/1.4.2_scheme2.png}
    \caption{Расположение грузов и призм на маятнике} 
\end{figure} }\\
Тогда положения грузов оказываются взаимно однозначно связаны соотношением:
 \begin{equation}Ml_1 = \frac{Lm_{ст}}{2} + m_{пр2}L + m_1b_1 +m2(b_2+L), \end{equation}
где $M=m_\text{ст}+m_{пр1}+m_{пр2}+m_1+m_2$ - полная масса маятника, $l_1$ - расстояние от острия призмы П1 до центра масс маятника. Здесь мы считаем, что призмы расположены на стержне симметрично и центр масс стержня равноудален на расстояние $\frac{L}{2}$ от призм. \\ \\
Тогда относительно подвеса
момент инерции тонкого стержня длиной $l_{ст}$ с призмами:
 \begin{center}$J_\text{ст}= m_{ст}(\frac{l_{ст}^2}{12}+(\frac{L}{2})^2) + m_{пр2}L^2$\end{center}
Момент инерции грузов на стержне:
 \begin{center}$J_\text{гр}= m_1(L-b_1)^2+m_2b_2^2$\end{center}
Суммарный момент инерции всего маятника:
 \begin{equation}J_\text{п}= MLl_2= J_\text{ст}+J_\text{гр}\end{equation}
Расчёт проводится по следующему алгоритму:
\begin{enumerate}
\item Выбираем и фиксируем некоторое значение расстояния $l_2$, удовлетворяющее условию (8).
\item По известным массам рассчитываем $J_\text{п}$ и $J_\text{cт}$, которые зависят только от величин $l_2$ и $L$.
\item Варьируя длину $b_2$ в пределах от $b_{2min}$ = 0 до $b_{2max} = \frac{l-L}{2}$, вычисляем по формуле (7) соответствующее положение первого груза.
$b_1$, а затем момент инерции грузов $J_{гр}$.
\item Строим график$J_{гр}(b_2)$и по точке пересечения с уровнем $J_{гр}=J_{п}-J_{ст}$ определяем искомые положения грузов $b_2$ и $b_1$.
\item Если расчётные положения грузов окажутся физически не реализуемы
(грузы слишком близко к призмам и т.п.), то следует выбрать другую
длину $l_2$ и повторить расчёт.
\end{enumerate}
\section{Оборудование}
\subsection{Используемое оборудование}
Оборотный маятник с двумя подвесными призмами и двумя грузами; электронный счётчик времени и числа колебаний; подставка с острием для определения положения центра масс маятника; закреплённая на стене консоль для подвешивания маятника; металлические линейки, штангенциркуль длиной 1 м.
\subsection{Инструментальные погрешности}
\vspace{0.15 cm}

\setdescription{leftmargin=0pt}
\begin{description}
    \item[штангенциркуль:] $\Delta_\text{шт}=\pm0,1$ мм.
    \item[весы:] $\Delta_\text{в}=\pm0,1$ г.
    \item[секундомер:] $\Delta_\text{с}=\pm0,01$ с.
\end{description}
\newpage
\section{Результаты измерений и обработка данных}

\begin{table}[!h]
\centering
\caption{Массы элементов}
\begin{tabular}{|c|c|}      
\hline
Величина & Значение \\  \hline
Масса стержня, $m_c$ & $1818,8 \pm 0,1$ г \\ \hline
Масса призмы 1, $m_\text{п1}$ & $75,6 \pm 0,1$ г\\ \hline
Масса призмы 2, $m_\text{п2}$ & $59,8 \pm 0,1$ г\\ \hline
Масса груза 1, $m_\text{г1}$ & $1493,8 \pm 0,1$ г\\ \hline
Масса груза 2, $m_{г2}$ & $1484,1 \pm 0,1$ г\\ \hline
Общая масса, $M$ & $4132,1 \pm 0,5$ г\\ \hline
\end{tabular}
\end{table}

\begin{table}[!h]
\centering
\caption{Расстояния (обозначения по рис. 2)}
\begin{tabular}{|c|c|}        \hline
Величина & Значение\\  \hline
Длина стержня, $l_c$ & $1000,8\pm 0,1$ мм\\ \hline
Расстояния между призмами 1, $L$ & $600,8\pm 0,3$ мм\\ \hline
Расчетное расстояние,$ b_1$ & $23,7$ см\\ \hline
Расчетное расстояние, $b_2$ & $12,5$ см\\ \hline
Расстояние от центра масс до П1,$l_1$ & $43,5\pm 0,1$ см\\ \hline
Расстояние от центра масс до П2, $l_2$ & $16,5\pm 0,1$ см\\ \hline
\end{tabular}
\end{table}

\noindent Калибровка положения груза 2:
\begin{table}[!h]
\centering
\begin{tabular}{|c|c|c|c|c|}      
\hline
№ & N, об & t2, c & t1, c & $\Delta$T, c \\  \hline
1 & 20 & 31,65 & 32,17 & 1,607 \\ \hline
2 & 20 & 31,53 & 31,99 & 1,589 \\ \hline
3 & 20 & 31,12 & 31,16 & 0,002\\ \hline
\end{tabular}
\end{table}

\noindent Более точное измерение периодов колебаний маятника:
\begin{table}[!h]
\centering
\begin{tabular}{|c|c|c|c|}        \hline
Т. подвеса & N, об & t, c & T, c\\  \hline
Призма 2 & 300 & $466,8 \pm 0,01\ c$ & $1,556$ \\ \hline
Призма 1 & 300 & $467,4 \pm 0,01\ c$ & $1,558$ \\ \hline
\end{tabular}
\end{table}

\noindent Конечные результаты для расчета $g$:
\begin{table}[!h]
\centering
\begin{tabular}{|c|c|}        \hline
Величина & Значение\\  \hline
$l_1$ & $43,5\pm 0,1$ см\\ \hline
$l_2$ & $16,5\pm 0,1$ см\\ \hline
$T_1$ & $1,556$ с\\ \hline
$T_2$ & $1,558$ с\\ \hline
$L$ & $60,08\pm 0,2$ см\\ \hline
\end{tabular}
\end{table} \\
\noindent
Рассчитаем начальную скорость пули по формуле:

\[g=(2\pi)^2\frac{l_1^2-l_2^2}{T_1^2l_1-T_2^2l_2}.\]
Рассчитаем погрешности:
\[
\frac{\sigma_{g_0}}{g_0} = \sqrt{ \left( \frac{\sigma_L}{L} \right)^2 + 4 \left( \frac{\sigma_T}{T} \right)^2 }. 
\]
Согласно (5):
\[\Delta g \approx \frac{2 l_2}{l_1 - l_2} \, \frac{\Delta T}{T} \, g_0 . \]
\noindent
Основной вклад в относительную погрешность $\Delta g$ вносят величины $\Delta T$ и $\Delta l = l_1-l_2$, поскольку являются разностями двух близких величин. Поэтому
\[
\frac{\sigma_{\Delta g}}{g} \approx \frac{2\beta \Delta T}{T} \sqrt{\left(\frac{\sigma_{\Delta T}}{\Delta T}\right)^2 + \left(\frac{\sigma_{\Delta l}}{\Delta l}\right)^2} = \sqrt{8\left(\beta \frac{\sigma_T}{T}\right)^2 + 8\left(\beta \frac{\Delta T}{T} \frac{\sigma_l}{\Delta l}\right)^2}.
\]
Тогда для полной относительной погрешности:
\[
\frac{\sigma_{g}}{g} \approx \sqrt{ \left( \frac{\sigma_{L}}{L} \right)^{2} + 4 \left( \frac{\sigma_{T}}{T} \right)^{2} + 8 \left( \beta \, \frac{\sigma_{T}}{T} \right)^{2} + 8 \left( \beta \, \frac{\Delta T}{T} \, \frac{\sigma_{l}}{\Delta l} \right)^{2} } \ .
\]	

\noindent $ g = 9,799 \pm 0,033$ м/$c^2$.

\section{Выводы}
В результате работы было измерено ускорение свободного падения при помощи оборотного маятника.

\end{document}
